\begin{statementbox}
  \textbf{\textcolor{CStatement}{这条结论解决什么问题？}}

  \medskip
  当分子和分母都含有 $x^2$ 时，直接观察分式往往看不出 $y$ 最小能是多少、最大能是多少。
  这时可以把问题\textbf{反过来问}：\textbf{先指定一个 $y$，有没有实数 $x$ 能使函数恰好等于它？}
  有，这个 $y$ 就能取到；没有，这个 $y$ 就不能取到。

  \medskip
  \begin{center}\begin{tikzpicture}[font=\small,>=stealth,line join=round]
\draw[gray!40,rounded corners=3pt] (0,0) rectangle (14.2,9.6);
\node at (7.1,9.0) {水平线有交点，才说明这个 $y$ 能取到};
\pgfmathsetmacro{\LL}{(3-sqrt(5))/2}
\pgfmathsetmacro{\UU}{(3+sqrt(5))/2}
\pgfmathsetmacro{\rr}{(sqrt(5)-1)/2}
\pgfmathsetmacro{\ss}{-(sqrt(5)+1)/2}
\begin{scope}[shift={(4.3,2.0)},x=.85cm,y=1.45cm]
\fill[blue!6] (-3.8,\LL) rectangle(3.0,\UU);
\draw[->,gray] (-4,0)--(3.3,0) node[right]{$x$};
\draw[->,gray] (0,-.25)--(0,3.5) node[above]{$y$};
\draw[teal!80!black,very thick,domain=-3.8:3.0,samples=200] plot(\x,{(\x*\x+1)/(\x*\x+2*\x+2)});
\draw[dashed,blue] (-3.8,\LL)--(3,\LL);
\draw[dashed,blue] (-3.8,\UU)--(3,\UU);
\draw[orange!85!black,thick] (-3.8,3)--(3,3);
\fill[blue] (\rr,\LL) circle(2pt);\fill[blue] (\ss,\UU) circle(2pt);
\node[left=5pt] at (-3.8,\LL){$L$};\node[left=5pt] at (-3.8,\UU){$U$};
\node[left=5pt,orange!85!black] at (-3.8,3){$3$};
\node[below left=4pt] at (0,0){$O$};
\end{scope}
\node[align=left,anchor=west] at (8.2,5.3){$F(x)=\dfrac{x^2+1}{x^2+2x+2}$\\[12pt]$L=\dfrac{3-\sqrt5}{2}$\\[10pt]$U=\dfrac{3+\sqrt5}{2}$\\[12pt]蓝线接触曲线，端点可取\\[7pt]橙线 $y=3$ 没有交点};
\node[text width=13cm,align=center] at (7.1,.6) {两条蓝色边界均可取到，值域为 $[L,U]$。};
\end{tikzpicture}
\end{center}
\end{statementbox}
\clearpage
\begin{statementbox}
\textbf{\textcolor{CStatement}{适用条件}}
  \begin{itemize}[leftmargin=*,itemsep=0.3em]
    \item 对于函数
    \[
      y=\frac{ax^2+bx+c}{dx^2+ex+f},\qquad d\ne0,
    \]
    其中 $a,b,c,d,e,f$ 是固定的实数，分子也允许没有 $x^2$ 项。
    \item 先看分母：若 $e^2-4df<0$，分母对任何实数 $x$ 都不为 $0$，所以 $x$ 可以取任意实数。本结论先讨论这一最省心的情况。
  \end{itemize}

  \medskip
  \textbf{\textcolor{CStatement}{核心结论}}
  把上式的分母乘过去并移项，得到
  \[
    \boxed{(a-dy)x^2+(b-ey)x+(c-fy)=0.}
  \]
  \begin{enumerate}[leftmargin=*,itemsep=0.4em]
    \item 当 $a-dy\ne0$ 时，它是真正的二次方程。某个 $y$ 能取到，\textbf{恰好等于}下面的式子不小于 $0$：
    \[
      \boxed{\Delta(y)=(b-ey)^2-4(a-dy)(c-fy)\ge0.}
    \]
    \item 当 $a-dy=0$，即 $y=\dfrac ad$ 时，$x^2$ 消失了。此时\textbf{不能只看 $\Delta(y)$}，必须把这个 $y$ 代回方程，检查是否真的能找到 $x$。
  \end{enumerate}

  \medskip
  \textbf{\textcolor{CStatement}{最后一步不能省}}
  解出 $\Delta(y)\ge0$ 给出的范围后，只保留其中 $a-dy\ne0$ 的值；再单独决定 $y=\dfrac ad$ 是否能够取到。两部分合起来才是完整值域。

  % ==== 新增：两条几何理解思路（保留原结论） ====
  \medskip
  \textbf{还有两种可以看到图形的理解方式：}
  把分母当成点的横坐标、分子当成纵坐标，$y$ 就是原点连线的斜率；
  对某些分母，还可以用平方差与勾股定理，把问题放到单位圆上讨论。
  这两种方法主要帮助理解\textbf{为什么会有最大、最小值}，并不能代替原来的分母检查。

\medskip
  \textbf{七种思路，解决的是同一个问题：}
  1. 判别式法；2. 两条抛物线比较法；3. 距离平方比法；
  4. 抛物线斜率法；5. 单位圆法；6. 倒数对称法；7. 导数法。
  \textbf{不是七种方法都适合每一道题：}1.可用于上面的通用形式；
  2.、4.主要讲清判别式与图像相切的联系；3.需要分子、分母能表示为距离平方；
  5.需要先找到平方和关系；6.须有特别的对称结构；7.适合已学导数的同学。
\end{statementbox}
